% --- To be compiled with XeLaTeX ---
% ---      Encoding: UTF-8        ---

\documentclass[a4paper, oneside, 11pt]{article}

%fontspec package provides a configurable interface for font selection, and allows complex font choices to be named and later reused. It's needed for XeLaTeX
\usepackage[cm-default]{fontspec}

% Unicode support
\usepackage{xunicode}
\usepackage{xltxtra}

% Default words and phrases in Greek (e.g. 'Περίληψη' instead of 'Abstract'). Also contains hyphenation rules for Greek Language
\usepackage{xgreek}

% Mathematical fonts, theorems etc.
\usepackage{amsfonts}
\usepackage{amsmath}
\usepackage{amsthm}

% Default page layout for consuming a larger portion of the page.
\usepackage{fullpage}

% Greek fonts (Computer Modern)
\setmainfont[Mapping=tex-text]{CMU Serif}

% Auxiliary commands
\newcommand{\HRule}{\rule{\linewidth}{0.5mm}}

\newtheorem{thm}{Θεώρημα}
\newtheorem{lm}[thm]{Λήμμα}

\begin{document}

\begin{titlepage}
\begin{center}

\includegraphics[width=0.3\textwidth]{/home/makis/.vim/skel/pyrforos.png}\\[1cm]

\textsc{\LARGE Σχολή Ηλεκτρολόγων Μηχανικών και Μηχανικών Υπολογιστών}\\[1.5cm]

\HRule \\[0.4cm]
{\huge \bfseries Γραφοθεωρία\\
\LARGE Ομάδα Ασκήσεων No. 1}\\[0.4cm]

\HRule \\[1.5cm]

\begin{center}
\textbf{Ομάδα 7}\\
Αξιώτης Κυριάκος\\
Αρσένης Γεράσιμος
\end{center}

\vfill

{\large \today}
\end{center}

\end{titlepage}


\section{Βαθμοί, κύκλοι, μονοπάτια}

\begin{enumerate}
\item[1.7 ($\star$)]
   Δείξτε ότι αν για κάποιο γράφημα $G$ ισχύει ότι
   $\textbf{διάμετρος}(G) \geq 2$, τότε το ίδιο θα ισχύει και για
   το απόκεντρο του G.

   \begin{proof}
   Έστω $\textbf{αποκ}(G)$ το σύνολο των κορυφών του $G$ που ανήκουν στο
   απόκεντρο και $H = G_{\textbf{αποκ}(G)}$ το εναγόμενο από το απόκεντρο
   υπογράφημα.

   Έστω ότι ήταν $\textbf{διάμετρος}(H) \leq 1$. Επειδή το απόκεντρο
   πρέπει να περιέχει τουλάχιστον 2 κορυφές, η διάμετρος δεν μπορεί
   να είναι 0 άρα έχουμε ότι
   $\textbf{διάμετρος}(H) = 1$, δηλαδή το $H$ είναι πλήρες γράφημα.

   Τότε όμως η διάμετρος του $G$ δεν μπορεί να είναι μεγαλύτερη ή ίση
   του 2 γιατί όλες οι αντιδιαμετρικές κορυφές βρίσκονται στο
   απόκεντρο και το απόκεντρο είναι πλήρες.

   Συνεπώς αν $\textbf{διάμετρος}(G) \geq 2$ τότε και
   $\textbf{διάμετρος}(H) \geq 2$.
   \end{proof}

\item[1.8 ($\star$)]
   Προσδιορίστε τη μέση απόσταση δύο κορυφών του γραφήματος $Q_r$ (δηλ.
   το μέσο όρο των αποστάσεων για όλα τα δυνατά ζεύγη διακεκριμένων κορυφών).

   \begin{proof}
      Ως γνωστόν οι κορυφές του υπερκύβου μπορούν να αριθμηθούν
      με δυαδικές συμβολοσειρές μήκους $r$. Δύο κορυφές συνδέονται
      με ακμή ανν οι συμβολοσειρές τους διαφέρουν μόνο σε μία θέση.

      Έστω μία κορυφή $x$. Το πλήθος των κορυφών που βρίσκονται σε
      απόσταση $d$ είναι ίσο με το πλήθος των κορυφών που οι συμβολοσειρές τους
      διαφέρουν σε $d$ ακριβώς θέσεις σε σχέση με την $x$. Δηλαδή υπάρχουν
      ${r \choose d}$ κορυφές σε απόσταση $d$.

      Συνεπώς έχουμε:

      \begin{align*}
         E[d] &= \frac{1}{{n(G) \choose 2}} \sum_{u, v \in V(G): u \neq v}
                  d(u, v)\\
              &= \frac{1}{{2^r \choose 2}} \sum_{u \in V(G)}
                  \sum_{v \in V(G): v \neq u} d(u, v)\\
              &= \frac{1}{\frac{2^r \cdot (2^r-1)}{2}}
                  \sum_{u \in V(G)} \sum_{k=1}^{r} k \cdot {r \choose k}\\
              &= \frac{2}{2^r (2^r - 1)} n(G) \sum_{k=1}^{r}
                  r {r-1 \choose k-1}\\
              &= \frac{2 \cdot 2^r \cdot r}{2^r (2^r - 1)} \sum_{k=0}^{r-1}
                  {r-1 \choose k}\\
              &= \frac{2 \cdot r \cdot 2^{r-1}}{2^r - 1}\\
              &= \frac{r \cdot 2^r}{2^r - 1}
      \end{align*}

   \end{proof}

\item[1.9 ($\star$)]
   Για κάθε θετικό ακέραιο $\alpha$ και για κάθε γράφημα $G$, το $V(G)$
   περιέχει περισσότερες από $\left( 1-\frac1\alpha \right) \cdot n(G)$
   κορυφές βαθμού αυστηρά μικρότερου του $2\alpha\delta^*(G)$.

   \begin{proof}
      Σύμφωνα με τον ορισμό έχουμε ότι
      $\delta^*(G) = \max \{ k\ |\ \exists H \subseteq G \text{ με }
       \delta(H) \geq k \}$.

      Θέλουμε να δείξουμε ότι:

      \[
         | \{ u\ |\ u \in V(G) \land d(u) < 2\alpha\delta^*(G) \} |
            > \left( 1 - \frac1\alpha \right) n(G)\\
      \]

      Έστω λοιπόν, προς απαγωγή σε άτοπο ότι:
      \begin{align*}
         | \{ u\ |\ u \in V(G) \land d(u) < 2\alpha\delta^*(G) \} |
            &\leq \left( 1 - \frac1\alpha \right) n(G)\\
         \Leftrightarrow
         n(G) - | \{ u\ |\ u \in V(G) \land d(u) \geq 2\alpha\delta^*(G) \} |
            &\leq \left( 1 - \frac1\alpha \right) n(G)\\ 
         \Leftrightarrow
         | \{ u\ |\ u \in V(G) \land d(u) \geq 2\alpha\delta^*(G) \} |
            &\geq \frac1\alpha n(G)\\
      \end{align*}

      Ισχύει όμως:

      \[ 2m(G) = \sum_{u \in V(G)} d(u) \geq
         \sum_{u \in V(G): d(u) \geq 2\alpha\delta*(G)} d(u)
         \geq \frac1\alpha n(G) \cdot 2\alpha\delta*(G)
         = 2n(G)\delta^*(G) \]

      \[ m(G) \geq n(G) \cdot \delta^*(G) \Leftrightarrow
         \epsilon(G) \geq \delta^*(G) \]

      Από το Πόρισμα 3.1 των σημειώσεων του μαθήματος γνωρίζουμε ότι
      $\delta^*(G) \geq \epsilon(G)$, συνεπώς θα έχουμε:

      \[ \epsilon(G) = \delta^*(G) \]

      TODO: Εδώ καταλήγουμε σε άτοπο γιατί αν ισχύει ισότητα τότε οι
      παραπάνω ανισότητες είναι αυστηρές και αυτό δεν μπορεί να ισχύει
      $[ \ldots ]$

   \end{proof}

\item[1.10 ($\star\star$)]
   Κάθε γράφημα $G$ με τουλάχιστον 2 κορυφές και
   $\epsilon(G) \geq 2$, έχει περιφέρεια το πολύ $2 \cdot \log_2 (n)$.

   \begin{proof}
      μπλα μπλα..
   \end{proof}
\end{enumerate}

\section{Άκυκλα γραφήματα}

\begin{enumerate}
   \item[2.10 $(\star)$]
      Σε κάθε δέντρο με $n$ κορυφές και διάμετρο τουλάχιστον
      $2k - 3$ υπάρχουν τουλάχιστον $n-k$ διαφορετικά μονοπάτια
      μήκους $k$.

      \begin{proof}

      *** Χρειάζεται Σχήμα (και ίσως αναδιατύπωση..) ***

      Έστω $u, v$ δύο αντιδιαμετρικοί κόμβοι με $d(u, v) = \text{diam}(u, v)$
      και έστω $P$ το μονοπάτι που τους ενώνει.
      Ονομάζουμε $w$ την κορυφή πάνω στο $P$ που απέχει $d(u, w) = k-1$
      από την $u$. Τέτοια κορυφή υπάρχει αφού $|P| = d(u, v) \geq 2k-3 > k-1$.

      %Υπάρχει κόμβος $w$ πάνω στο
      %$P$ τέτοιος ώστε είτε $d(u, w) \geq k-1$ είτε $d(w, v) \geq k-1$
      %(διαφορετικά θα είχαμε $\text{diam}(G) = d(u, v) = d(u, w) + d(w, v) \leq
      %2(k-2) = 2k - 4$).

      %Υποθέτουμε λοιπόν χωρίς βλάβη της γενικότητας ότι $d(u, w) \geq k-1$
      %και μάλιστα επιλέγουμε το κοντινότερο τέτοιο $w$ στο $u$, δηλαδή
      %$d(u, w) = k-1$.
      
      Στο μονοπάτι $P'$ από $u$ στον $w$ υπάρχουν ακριβώς $k$ κορυφές.
      Θα δείξουμε ότι με αφετηρία κάθε μία από τις υπόλοιπες $n-k$ κορυφές
      μπορούμε να δημιουργήσουμε διαφορετικά μονοπάτια μήκους $k$.

      Έστω μια κορυφή $x$ που δεν ανήκει στο $P'$. Θα δημιουργήσουμε ένα
      μονοπάτι $T_x$ μήκους $k$ διακρίνοντας δύο περιπτώσεις:

      \begin{enumerate}
         \item Αν $w \in P(x, u)$ όπου $P(x, u)$ το μονοπάτι από $x$ σε
               $w$ στο δέντρο τότε θέτουμε $T_x$ το πρόθεμα μήκους $k$ του
               μονπατιού (δηλαδή το $T_x$ περιλαμβάνει την αφετηρία $x$ και
               τους επόμενους $k-1$ κόμβους).

               Τέτοιο πρόθεμα υπάρχει πάντα γιατί το $P(x, u) \geq
               P(x, w) + 1 = k$.

         \item Αν $w \notin P(x, u)$ τότε θεωρούμε το μονοπάτι $P(x, v)$
               για το οποίο ισχύει $w \in P(x, v)$ και θέτουμε $T_x$ το
               πρόθεμα μήκους $k$ αυτού του μονοπατιού.

               Όπως πριν, εξασφαλίζουμε την ύπαρξη τέτοιου προθέματος
               από το γεγονός ότι $P(x, v) \geq P(w, v) + 2 \geq
               (2k-3) - (k-1) + 2 = k$
      \end{enumerate}

      Τα παρακάτω λήμματα μας εξασφαλίζουν ότι τα μονοπάτια που προκύπτουν
      από την παραπάνω διαδικασία είναι όλα διαφορετικά μεταξύ τους.

      \begin{lm}
         Έστω δύο μονοπάτια $P_1, P_2$ σε ένα δέντρο που έχουν προκύψει ώς
         πρόθεμα (προσανατολισμένων) μονοπατιών από την κορυφή $x_1$ στην $u$ και
         από την $x_2$ στην $u$ αντίστοιχα όπου οι $x_1, x_2, u$ διαφορετικές
         μεταξύ τους κορυφές. Τότε $P_1 \neq P_2$\footnote{Ορίσαμε τα $P_1, P_2$
         ώς πρόθεμα προσανατολισμένων μονοπατιών, δηλαδή μονοπατιών
         με συγκεκριμένη αφετηρία και πέρας όμως από τη στιγμή που τα ορίζουμε
         τα θεωρούμε πλέον μη-προσανατολισμένα και έτσι έχει νόημα η σύγκριση
         $P_1 \neq P_2$.}
      \end{lm}
      \begin{proof}
         Από τον ορισμό των $P_1, P_2$ βλέπουμε ότι ο μόνος τρόπος να
         είναι το ίδιο μονοπάτι είναι αν έχουν ώς άκρα τις κορυφές
         $x_1, x_2$.

         Αυτό σημαίνει ότι $x_2 \in P(x_1, u), x_1 \in P(x_2, u)$ το οποίο
         είναι άτοπο άρα $P_1 \neq P_2$.
      \end{proof}

      \begin{lm}
      Έστω δύο μονοπάτια $T_{x_1}, T_{x_2}$ για $x_1 \neq x_2$ που έχουν προκύψει
      από τις περιπτώσεις (α'), (β') αντίστοιχα. Τότε $T_{x_1} \neq T_{x_2}$.
      \end{lm}

      \begin{proof}
      Έχουμε ότι
      $x_2 \notin P(x_1, u)$ γιατί διαφορετικά είτε θα είχαμε
      $x_2 \in P(x_1, w)$ και τότε η $x_2$ θα ήταν στην περίπτωση (α') είτε
      $x_2 \in P(w, u) = P'$ το οποίο
      δεν μπορεί να συμβαίνει αφού οι κορυφές του $P'$ δεν είναι αφετηρίες
      μονοπατιών.

      Συνεπώς, το $T_{x_1}$ που είναι υποσύνολο του $P(x_1, u)$ δεν μπορεί
      να περιέχει την $x_2$, άρα τα $T_{x_1}, T_{x_2}$ έχουν τουλάχιστον
      μία κορυφή διαφορετική και έτσι είναι διαφορετικά.
      \end{proof}

      Σε κάθε περίπτωση λοιπόν τα $n-k$ μονοπάτια που δημιουργήσαμε
      είναι όλα διαφορετικά μεταξύ τους.
      \end{proof}
\end{enumerate}

\section{Συνεκτικότητα}
\begin{enumerate}
   \item[3.9 $(\star)$]
      Ένα γράφημα είναι δισυνεκτικό αν και μόνο αν μπορεί να κατασκευαστεί
      αρχίζοντας από το $K_3$ και εφαρμόζοντας μία ακολουθία μετασχηματισμών
      που μπορεί να είναι,

      \begin{itemize}
         \item Υποδιαίρεση ακμής.
         \item Πρόσθεση ακμής.
      \end{itemize}

      \begin{proof}
         Θα δούμε τις δύο κατευθύσεις του θεωρήματος ξεχωριστά.
         \begin{itemize}
            \item $\Leftarrow$
               Το $K_3$ είναι δισυνεκτικό άρα θα πρέπει να δείξουμε ότι
               οι παραπάνω μετασχηματισμοί διατηρούν αναλλοίωτη τη
               συνεκτικότητα.

               Πράγματι:

               \begin{itemize}
               \item
               Με την προσθήκη ακμής όλα τα μονοπάτια που υπάρχαν στο
               αρχικό γράφημα διατηρούνται. Έτσι, όσα εσωτερικώς διακεκριμένα
               μονοπάτια υπήρχαν μεταξύ ζευγών κορυφών συνεχίζουν να
               υπάρχουν και έτσι από το Θεώρημα Menger έχουμε ότι το
               γράφημα θα συνεχίσει να είναι δισυνεκτικό.

               \item
               Για την υποδιαίρεση ακμής, θα πρέπει να βεβαιωθούμε ότι
               δεν μπορούμε να αποσυνδέσουμε το γράφημα με την αφαίρεση
               μία κορυφής και συγκεκριμένα της κορυφής που βάλαμε με
               την υποδιαίρεση.

               Η αφαίρεση της νέας κορυφής ισοδυναμεί με αφαίρεση της
               υποδιαιρούμενης ακμής στο αρχικό γράφημα. Έστω $\{u, v\}$
               αυτή η ακμή. Επειδή το γράφημα
               αρχικά ήταν δισυνεκτικό, θα υπάρχει τουλάχιστον άλλο ένα
               μονοπάτι από την $u$ προς την $v$ άρα το γράφημα παραμένει
               συνεκτικό και μετά την αφαίρεση της $\{u, v\}$.
               \end{itemize}

            \item $\Rightarrow$
               Θα δείξουμε ότι αν ένα γράφημα $G$ είναι δισυνεκτικό τότε
               είτε:
               \begin{enumerate}
               \item[(1)] θα είναι το $K_3$, είτε
               \item[(2)] θα περιέχει μια κορυφή
               βαθμού 2 της οποίας οι γείτονες να μην είναι συνδεδεμένοι
               απευθείας
               (θα καλούμε τέτοιες κορυφές μή-απλοϊδείς) και η διάλυσή
               της δημιουργεί δισυνεκτικό γράφημα, είτε
               \item[(3)] θα περιέχει
               μία ακμή της οποίας η αφαίρεση οδηγεί σε δισυνεκτικό
               γράφημα.
               \end{enumerate}

               Έτσι, για κάθε δισυνεκτικό γράφημα μπορούμε να εφαρμόσουμε
               μία ακολουθία από διαλύσεις κορυφών και αφαιρέσεις ακμών
               μέχρι να καταλήξουμε στο $K_3$ και η αντίστροφη διαδικασία
               είναι που μας παράγει το $G$ από το $K_3$ όπως ζητάει η
               εκφώνηση.

               Αν το γράφημα $G$ περιέχει μια ακμή της οποίας η αφαίρεση
               διατηρεί το γράφημα δισυνεκτικό τότε έχουμε τελειώσει γιατί
               ισχύει το (3). Επομένως αρκεί να εξετάσουμε την περίπτωση
               όπου για όλες τις ακμές $e \in E(G)$ ισχύει
               $\kappa(G \backslash e) < 2$.

               Παρατηρούμε ότι $\kappa(G) = 2$ γιατί διαφορετικά, έστω
               ότι $\kappa(G) \geq 3$ τότε με την αφαίρεση μιας ακμής
               η συνεκτικότητα δεν θα έπρεπε να πέφτει πάνω από μία
               μονάδα (Παρατήρηση 5.7 των σημειώσεων του μαθήματος),
               όμως υποθέσαμε ότι η αφαίρεση οποιαδήποτε ακμής
               οδηγεί σε συνεκτικότητα μικρότερη του 2, δηλαδή έχουμε
               μείωση της συνεκτικότητας κατά 2 που είναι άτοπο.

               Σύμφωνα με το Θεώρημα Halin (συγκεκριμένα με το
               αντιθετο-αντίστροφό του) έχουμε ότι $\delta(G) \leq \kappa(G)
               = 2$. Ο ελάχιστος βαθμός ενός δισυνεκτικού γραφήματος
               δεν μπορεί να είναι μικρότερος του 2, άρα $\delta(G) = 2$.
               Έστω λοιπόν $u$ μια κορυφή βαθμού 2 και έστω $x, y$ οι
               γείτονές τις.

               Έστω τώρα ότι $\{x, y\} \in E(G)$. Αν το γράφημα έχει μόνο
               3 κορυφές τότε είναι το $K_3$ και έχουμε τελειώσει.
               Διαφορετικά έστω ότι έχει τουλάχιστον άλλη μία κορυφή $w$
               η οποία συνδέεται στην $x$. Επειδή το γράφημα είναι δισυνεκτικό
               θα πρέπει η αφαίρεση της $x$ να μην το αποσυνδέει, συνεπώς
               θα πρέπει να υπάρχει μονοπάτι $P$ από την $w$ στην
               $y$ που να μην χρησιμοποιεί την κορυφή $x$. Tότε όμως
               μεταξύ της $x$ και της $y$ θα υπήραν 3 εσωτερικά διακεκριμένα
               μονοπάτια: 1) $(x, y)$, 2) $(x, u, y)$, 3) $P$. Άτοπο
               γιατί τώρα η αφαίρεση της $\{x, y\}$ διατηρεί το γράφημα
               δισυνεκτικό.\footnote{Υπάρχει περίπτωση το $P$ να χρησιμοποιεί
               την $u$ ώς ενδιάμεσο κόμβο. Σε αυτή την περίπτωση θεωρούμε
               το $P'$ από το $w$ στο $u$ και δείχνουμε ότι υπάρχουν
               3 εσωτερικά διακεκριμένα μονοπάτια από το $x$ στο $u$.}

               Άρα η $u$ είναι μη απλοϊδής κορυφή βαθμού 2 και μένει να
               δείξουμε ότι η διάλυσή της διατηρεί τη συνεκτικότητα. Αυτό
               προκύπτει από το θεώρημα Menger αφού ό,τι μονοπάτια υπήρχαν
               πριν μεταξύ κορυφών συνεχίζουν να υπάρχουν.
         \end{itemize}
      \end{proof}

   \item[3.10 $(\star\star)$]
Για κάθε $k$ κορυφές ενός $k$-συνεκτικού γραφήματος, υπάρχει κύκλος που να τις περιέχει όλες.
\begin{proof}
Έστω $k$ κορυφές του γραφήματος $G$ και $C$ κύκλος που περιέχει όσο το δυνατόν περισσότερες από τις $k$ κορυφές. Έστω $S$ το σύνολο των $k$ κορυφών. Αν $|C|=k$, τελειώσαμε. 
Διαφορετικά, έστω $|C|=l$ και $u$ μία από τις 
$k$ κορυφές, η οποία δεν 
βρίσκεται στον κύκλο. Από το Λήμμα 1, υπάρχουν $min(|C|,k)$ εσωτερικώς διακεκριμένα μονοπάτια από το $u$ προς τις κορυφές του κύκλου, και κανένα δεν τελειώνει στην ίδια κορυφή του κύκλου. Έστω $v_i$ μία
απαρίθμηση των κορυφών του κύκλου (με τη σειρά που εμφανίζονται πάνω στον κύκλο) οι οποίες αποτελούν άκρο κάποιου μονοπατιού από τα παραπάνω 
και $P_i$ τα αντίστοιχα μονοπάτια. Επίσης έστω $F_i$ το μονοπάτι από 
την $v_i$ στην $v_{i+1}$ το οποίο δεν περιέχει καμία άλλη από τις $v_j$. (Έχουμε θεωρήσει ότι $v_{min(|C|,k)+1}\equiv v_1$). Αν ο κύκλος έχει μήκος $l$, τότε περιέχει μόνο κορυφές από το $S$. Ο κύκλος
$v_1, P_1, u, P_2, v_2, v_3, ..., v_l, v_1$ περιέχει $l+1$ στοιχεία του $S$, άτοπο. Αν έχει μήκος $>l$, τότε οι κορυφές $v_i$ είναι τουλάχιστον $l+1$. Αυτό σημαίνει ότι υπάρχουν τουλάχιστον $l+1$
διαφορετικά μονοπάτια $F_i$. Άρα θα υπάρχει ένα $F_i$ το οποίο δεν περιέχει στο εσωτερικό του καμία κορυφή του $S$. Τότε, ο κύκλος $v_1, F_1, v_2, ..., v_i, P_i, u, P_{i+1}, v_{i+1}, F_{i+1}, ...,v_1$ 
έχει $l+1$ στοιχεία του $S$, άτοπο.
Άρα για κάθε σύνολο $k$ κορυφών, υπάρχει κύκλος που τις περιέχέι όλες.
\end{proof}

Λήμμα 1: Έστω $k$-συνεκτικό γράφημα, κύκλος του με τουλάχιστον $l$ κορυφές με $l<k$ και τυχαία κορυφή $u$ εκτός του κύκλου. Τότε υπάρχουν $l$ κορυφές του κύκλου $v_1, v_2, ..., v_l$ και 
εσωτερικώς διακεκριμένα μονοπάτια $P_i = u...v_i$ για κάθε $1\leq i\leq l$.
\begin{proof}
Έστω μία νέα κορυφή $v$ που συνδέεται με ακμή με όλες τις κορυφές του κύκλου. Δηλαδή θεωρούμε γράφημα $G$ με $V(G')=V(G)\cup \{v\}$ και $E(G')=E(G)\cup \{(v,x)|x\in C\}$. Το $G$ είναι $l$-συνεκτικό:
Αν σβήσουμε $l-1$ κορυφές και σε αυτές περιέχεται η $v$, τότε οι κορυφές που απομένουν συνδέονται λόγω της $k$-συνεκτικότητας του αρχικού γραφήματος. Σε διαφορετική περίπτωση, θα σβηστούν το πολύ $l-1$
κορυφές του κύκλου και συνεπώς θα μείνει τουλάχιστον μία άκμή από την $v$ προς μια κορυφή του κύκλου, άρα το γράφημα θα παράμείνει συνεκτικό. Αφού το γράφημα είναι $l$-συνεκτικό, θα υπάρχουν $l$ εσωτερικώς
διακεκριμένα μονοπάτια από την κορυφή $u$ στην κορυφή $v$. Κάθε ένα από αυτά τα μονοπάτια περνάει από τουλάχιστον μία κορυφή του κύκλου. Για κάθε μονοπάτι $P = u...v$, θεωρούμε την πρώτη φορά που περνάει
από μία κορυφή του κύκλου. Έστω ότι αυτή είναι η $x_i$. Το σύνολο των μονοπατιών $\{P_i = u...x_i\}$ είναι το ζητούμενο, αφού τα μονοπάτια είναι εσωτερικώς διακεκριμένα και καταλήγουν σε $l$ διαφορετικές
κορυφές του κύκλου.

\end{proof}

\end{enumerate}
\section{Εμβαπτίσεις}
\section{Δομές σε γραφήματα}

\begin{enumerate}
   \item[5.9 $(\star)$]
      Κάθε γράφημα περιέχει τουλάχιστον $\frac{m(G)(4m(G) - n^2(G))}{3n(G)}$
      τρίγωνα.

      \begin{proof}

      Έστω μια ακμή $\{u, v\}$. Η ιδέα είναι να βρούμε το ελάχιστο πλήθος
      τριγώνων στα οποία μπορεί να ανήκει αυτή η ακμή και έτσι μετά αθροίζοντας
      κατάλληλα να μπορέσουμε να φράξουμε από κάτω το συνολικό πλήθος των
      τριγώνων του γραφήματος.

      Ορίζουμε $U = N_G(u) \backslash v,
      V = N_G(v) \backslash u$. Ισχύει ότι $|U| + |V| = d(u) + d(v) - 2$.
      Επίσης, $|U \cup V| \leq n(G) - 2$ αφού δεν υπάρχουν πάνω από
      $n(G)-2$ κορυφές που να απομένουν στο γράφημα.

      Άρα, έχουμε:

      \[ |U \cap V| = |U| + |V| - |U \cup V| \geq d(u) + d(v) - n(G) \]

      Κάθε κορυφή που ανήκει στο $U \cap V$ δημιουργεί τρίγωνο με τις
      κορυφές $u, v$. Άρα το πλήθος των τριγώνων $|T_{\{u, v\}}|$ που μπορεί
      να ανήκει
      η ακμή $\{u, v\}$ είναι τουλάχιστον $d(u) + d(v) - n(G)$.

      Αν συμβολίσουμε με $T$ το σύνολο των τριγώνων του $G$ έχουμε:

      \[ 3|T| = \sum_{\{u, v\} \in E(G)} T_{\{u, v\}} \]

      επειδή κάθε τρίγωνο περιέχει 3 ακμές.

      Συνεπώς:

      \begin{align*}
         |T| &\geq \frac13 \sum_{\{u, v\} \in E(G)} ( d(u) + d(v) - n(G) )\\
             &= \frac13 \sum_{\{u, v\} \in E(G)} ( d(u) + d(v) ) -
                \frac{n(G)m(G)}3 \\
             &= \frac13 \sum_{u \in V(G)} d^2(u) - \frac{n(G)m(G)}3\\
             &\geq \frac1{3n(G)} \left( \sum_{u \in V(G)} d(u) \right)^2
                 - \frac{n(G)m(G)}3\\
             &= \frac{4m^2(G)}{3n(G)} - \frac{n(G)m(G)}3\\
             &= \frac{m(G)(4m(G) - n^2(G))}{3n(G)}\\
      \end{align*}

      Όπου το 4ο βήμα προκύπτει από την ανισότητα Cauchy-Schwarz:

      \[ d(u_1)\cdot1 + d(u_2)\cdot2 + \ldots + d(u_n)\cdot1
         \leq (d^2(u_1) + \ldots + d^2(u_n)) \cdot (1 + \ldots + 1)
         = (d^2(u_1) + \ldots + d^2(u_n)) \cdot n \]

      \end{proof}

   \item[5.10 $(\star \star)$]
THELEI FTIAKSIMO
Δείξτε ότι ένα πολυγράφημα είναι σειριακό-παράλληλο αν είναι 2-συνεκτικό και δεν περιέχει καμία υποδιαίρεση του $K_4$ ως ελάσσον. Ένα γράφημα καλείται σειριακό-παράλληλο αν μπορεί
να προκύψει από το $K_2$ μετά από σειρά υποδιαιρέσεων ακμών ή διπλασιασμών ακμών (δηλαδή αντικατάσταση μιας ακμής από μια διπλής πολλαπλότητας με τα ίδια άκρα).
	\begin{proof}
Αρχικά, αν ένα γράφημα περιέχει κάποια υποδιαίρεση του $K_4$ ως ελάσσον, τότε περιέχει και το $K_4$ ως ελάσσον, αφού η διάλυση κορυφής είναι η αντίστροφη πράξη της υποδιαίρεσης ακμής και γνωρίζουμε
ότι η σύνθλιψη ακμής μπορεί να προσομοιώσει την διάλυση κορυφής.
Θα δείξουμε κάτι πιο ισχυρό, δηλαδή ότι αν ένα πολυγράφημα είναι συνεκτικό και δεν περιέχει το $K_4$ ως ελάσσον, τότε είναι σειριακό-παράλληλο.
Για 2 κορυφές ισχύει, αφού έχουμε το $K_2$ που είναι σειριακό-παράλληλο.
Για 3 κορυφές επίσης ισχύει. Αν έχουμε το $P_3$, τότε προκύπτει από το $K_2$ με μία υποδιαίρεση ακμής. Αν έχουμε το $K_3$, αυτό μπορεί να προκύψει από την εξής ακολουθία κινήσεων: 
$K_2$->διπλασιασμός ακμής, υποδιαίρεση της μίας ακμής.
Θεωρούμε το γράφημα $G$ με τον ελάχιστο αριθμό κορυφών, το οποίο είναι συνεκτικό και δεν περιέχει το $K_4$ ως ελάσσον και δεν είναι σειριακό-παράλληλο. Από το Λήμμα 1, το γράφημα $G$ δεν μπορεί
να είναι 3-συνεκτικό. Αν το $G$ δεν είναι 2-συνεκτικό, τότε έχει κορυφή τομής. Έστω
Έστω ένας 2-διαχωριστής $u$, $v$ και $G'$ μία συνεκτική συνιστώσα που προκύπτει μετά τη διαγραφή των κορυφών $u$, $v$. Έστω γράφημα $H$ με 
$V(H)=V(G')\cup {u,v}$ και $E(H) = {(x,y) | x\in V(H), y\in V(H), (x,y)\in G, (x,y)\neq (u,v)}$. Το γράφημα $H$ είναι συνεκτικό, διότι σε διαφορετική περίπτωση κάποια από τις κορυφές 
$u$, $v$ θα αποτελούσε κορυφή τομής. Αν το $H$ δεν είναι 2-συνεκτικό, τότε λόγω του Λήμματος 2 είναι και σειριακό-παράλληλο με άκρα τα $u$, $v$. Αν είναι 2-συνεκτικό αλλά ότι 3-συνεκτικό, τότε 
λόγω του Λήμματος 3 
και της υπόθεσης ελαχιστότητας είναι σειριακό-παράλληλο με άκρα τα $u$, $v$. Λόγω του Λήμματος 3, το $H$ δεν μπορεί να είναι 3 συνεκτικό. Θεωρούμε λοιπόν το γράφημα $Κ_2$. Έστω $x$ το πλήθος των
συνεκτικών συνιστωσών που προκύπτουν στο $G$ μετά τη διαγραφή των $u$, $v$. Αν στο $G$ υπάρχει η ακμή $(u,v)$, τότε διπλασιάζουμε $x$ φορές την ακμή του $K_2$. Διαφορετικά τη διπλασιάζουμε $x-1$ φορές.
Τώρα, εφόσον έχουμε αποδείξει ότι κάθε συνιστώσα, μαζί με τις κορυφές $u$,$v$ είναι σειριακό-παράλληλο γράφημα με άκρα τα $u$, $v$, μπορεί να προκύψει με μια σειρά διπλασιασμών και υποδιαιρέσεων ακμών
από το $K_2$, δηλαδή από μία από τις ακμές που προέκυψαν από το διπλασιασμό. Αν επαναλάβουμε αυτή τη διαδικασία για κάθε συνιστώσα, θα προκύψει το γράφημα $G$. Αυτό σημαίνει ότι το $G$ είναι 
σειριακό-παράλληλο, το οποίο είναι άτοπο. Άρα η υπόθεση δεν ισχύει και κάθε 2-συνεκτικό γράφημα που δεν περιέχει το $K_4$ ως ελάσσον είναι σειριακό-παράλληλο. Παρατήρηση: Η συνθήκη ότι ένα γράφημα
δεν περιέχει το $K_4$ ως ελάσσον, με την προϋπόθεση ότι είναι συνεκτικό, είναι ικανή για την απόδειξη του ότι είναι σειριακό-παράλληλο.
	\end{proof}     
Λήμμα 1: Για κάθε γράφημα $G$ με $n(G)\geq 4$, ισχύει ότι $\kappa(G)\geq 3 \Rightarrow K_4 \subseteq_{\epsilon\lambda} G$.
	\begin{proof}
Είναι το Πόρισμα 5.44 από τις σημειώσεις του μαθήματος.
	\end{proof}
Λήμμα 2: Κάθε συνεκτικό γράφημα που δεν είναι 2-συνεκτικό είναι 
	\begin{proof}

	\end{proof}
 
\end{enumerate}

\section{Χρωματισμοί και άλλα}
\begin{enumerate}

\item[6.9 ($\star\star$)]
   Ένα γράφημα λέγεται άρτιο αν όλες οι κορυφές έχουν άρτιο βαθμό. Δείξτε ότι αν το G είναι συνεκτικό γράφημα, τότε 
   $|\{H\subseteq_{\text{πα}} G | H \text{είναι άρτιο}\}| = 2^{m(G)-n(G)+1}$.

   \begin{proof}
Θεωρούμε $S = \{H\subseteq_{\pi \alpha} G | H \text{άρτιο}\}$.
     Θα ορίσουμε μία 1-1 και επί συνάρτηση $f$ από το σύνολο 
$A = \{H | H\subseteq_{\pi\alpha} G\}$, δηλαδή το σύνολο των παραγόμενων γραφημάτων 
του $G$, στο $B = S\times \{X\subseteq V(G) | |X| mod 2 = 0\}$, 
δηλαδή το καρτεσιανό γινόμενο του συνόλου των άρτιων παραγόμενων γραφημάτων 
με την οικογένεια υποσυνόλων του $V(G)$ με άρτιο πληθάριθμο. Το σύνολο των 
παραγόμενων γραφημάτων του $G$ έχει πληθάριθμο $2^{m(G)}$, αφού κάθε ακμή 
μπορεί να υπάρχει ή να μην υπάρχει στο παραγόμενο υπογράφημα. 
Επίσης η οικογένεια υποσυνόλων του $V(G)$ με άρτιο πληθάριθμο έχει πληθάριθμο 
$2^{n(G)-1}$, αφού έχουμε 2 επιλογές για κάθε κορυφή (θα μπει ή δεν θα μπει 
στο υποσύνολο), εκτός από την τελευταία, της οποίας η τοποθέτηση καθορίζεται 
μοναδικά από το αν το υποσύνολο έχει άρτιο ή περιττό αριθμό κορυφών. Λόγω 
του λήμματος 1, η $f$ είναι 1-1 και επί, άρα έχουμε ότι 
$2^{m(G)} = |S| \cdot 2^{n(G)-1} \Rightarrow |S| = 2^{m(G)-n(G)+1}$, το οποίο 
είναι το ζητούμενο.
   \end{proof}


Λήμμα 1: Υπάρχει 1-1 και επί συνάρτηση από το σύνολο $A$ στο σύνολο $B$.
\begin{proof}
Για κάθε ζευγάρι κορυφών $i$, $j$ με $i\neq j$, ορίζουμε $P_{ij}$ ένα μονοπάτι μεταξύ τους στο $G$. Αυτό προφανώς υπάρχει, αφού το $G$ είναι συνεκτικό.
Ορισμός $f$: Έστω $Z\in A$ και $T$ το σύνολο των κορυφών του $Z$ με περιττό βαθμό. Είναι γνωστό ότι $|Z|\text{mod} 2=0$. Διαμερίζουμε τις κορυφές του $Z$ σε ζευγάρια
$(a_i, b_i)$ (με κάποιο μονοσήμαντο τρόπο, πχ αριθμούμε τις κορυφές του $Z$ $u_1,u_2,...,u_k$ και βάζουμε τα ζευγάρια $(u_1,u_2), ..., (u_{k-1},u_k)$)
και για κάθε ζευγάρι θεωρούμε το μονοπάτι $P_{a_i b_i}$. Για κάθε ακμή πάνω σε αυτό το μονοπάτι, αν υπάρχει στο $Z$ τότε την αφαιρούμε, ενώ αν δεν
υπάρχει την προσθέτουμε. Είναι εύκολο να δούμε ότι αυτός ο μετασχηματισμός διατηρεί την αρτιότητα των βαθμών των ενδιάμεσων κόμβων, και επίσης πλέον οι $a_i$, $b_i$
έχουν άρτιο βαθμό. Κάνοντας αυτό το μετασχηματισμό για κάθε ζευγάρι, θα καταλήξουμε με ένα άρτιο γράφημα $U$. Ορίζουμε $f(Z)=U\times T$. Ουσιαστικά η $f$ μετασχηματίζει ένα
γράφημα σε άρτιο, αλλά επιστρέφει και την πληροφορία του ποιοι κόμβοι ήταν περιττοί. Αντίστροφα, αν έχουμε ένα άρτιο γράφημα $U$ και ένα υποσύνολο $T$ του $V(G)$ με άρτιο πληθάριθμο,
θεωρούμε τη διαμέριση του $T$ σε ζευγάρια και για κάθε ζευγάρι εφαρμόζουμε τον ίδιο μετασχηματισμό που ορίσαμε παραπάνω. Έτσι θα πάρουμε ξανά το γράφημα $Z$ με $f(Z)=U\times T$.
Συνεπώς η $f$ είναι 1-1 και επί.
\end{proof}

\item[6.10 ($\star\star$)]
	Δείξτε ότι υπάρχει θετική σταθερά $c$, τέτοια ώστε αν για κάποιο γράφημα $G$ ισχύει ότι $\delta (G)\geq k$, τότε το $G$ περιέχει $c\cdot k^2$ ακμοδιακεκριμένους κύκλους.

\begin{proof}
Έστω $\delta(G)\geq k\geq 4$. Λόγω του λήμματος 2, έχουμε $\geq \lfloor \frac{k-1}{3}\rfloor$ εσωτερικώς διακεκριμένους κύκλους. Διαγράφουμε τις ακμές όλων αυτών των κύκλων. Στο γράφημα
$G'$ που θα προκύψει έχουμε $\delta(G)\geq k-2$. Εφαρμόζουμε επαναληπτικά την ίδια διαδικασία, έως ότου το γράφημα που απόμένει έχει $\delta(G')<4$. Συνολικά αυτή η διαδικασία θα επαναληφθεί
τουλάχιστον $\lfloor\frac{k}{2}\rfloor -1$ φορές. Οι ακμοδιακεκριμένοι κύκλοι που θα έχουμε συνολικά λοιπόν θα είναι τουλάχιστον $\lfloor \frac{k-1}{3}\rfloor + \lfloor \frac{k-3}{3}\rfloor + ... + 1 + 0 = \Theta (k^2)$.
\end{proof}

Λήμμα 1: Αν $\delta(G)\geq 4$, υπάρχει κύκλος με μήκος $\leq 2\cdot log_2 n$.
\begin{proof}
Έχουμε $m\geq \frac{\delta(G)\cdot n}{2}\geq 2n$, άρα η πυκνότητα είναι τουλάχιστον 2. Αυτό που μένει έχει αποδειχθεί στην άσκηση 1.10.
\end{proof}

Λήμμα 2: Σε κάθε γράφημα $G$ με $\delta(G)\geq k\geq 4$ υπάρχουν τουλάχιστον $\lfloor \frac{k-1}{3}\rfloor$ εσωτερικώς διακεκριμένοι κύκλοι.
\begin{proof}
Έστω ένας ελάχιστος κύκλος $C$. Αυτος λόγω του λήμματος 1 θα έχει μήκος το πολύ $2\cdot log_2 n$. Επίσης καμία κορυφή $u\in G-C$ δεν μπορεί να έχει πάνω από 3 ακμές προς κορυφές του $G$.
Αν είχε, τότε έστω δύο από αυτές και οι αντίστοιχες κορυφές του κύκλου. Αυτές θα είχαν απόσταση $\leq \lfloor \frac{|C|}{2}\rfloor$ στον $C$, άρα χρησιμοποιώντας αυτές τις δύο ακμές, 
θα υπήρχε κύκλος με μέγεθος το πολύ
$\lfloor\frac{|C|}{2}\rfloor + 2$, το οποίο για $|C|\geq 5$ είναι άτοπο αφού δημιουργεί κύκλο μικρότερο από τον ελάχιστο. Για $|C|=3$, είναι προφανές ότι δεν μπορούμε να έχουμε πάνω από 3
ακμές από κάποια κορυφή προς τις κορυφές του $C$, ενώ για $|C|=4$ αν είχαμε 4 ακμές προς κορυφές το $C$, θα σχηματιζόταν κύκλος μήκους 3, άτοπο.
Από το παραπάνω συμπεραίνουμε ότι το εναγόμενο γράφημα $G'$ του $G$ με σύνολο κορυφών το $G-C$ θα έχει $\delta(G') \geq k-3$. Εφαρμόζοντας επαναληπτικά την ίδια διαδικασία στο εναγόμενο 
γράφημα, μέχρι ο ελάχιστος
βαθμός του αντίστοιχου εναγόμενου γραφήματος να γίνει μικρότερος από 4, έχουμε συνολικά τουλάχιστον $\lfloor \frac{k-1}{3}\rfloor$ εσωτερικώς διακεκριμένους κύκλους.
\end{proof}

\end{enumerate}

\end{document}
